% problem_id: S051-drops-D3
% source_id: S051 (DROPS/LIPIcs)
% source_file: t12.pdf
% statement_locator: printed p. PDF p. 16 = printed p. 12:16 (last item on the page; the object the statement names, iEF^{tt(h,n_0)}, is defined by Definition 17 on the same page = PDF p. 16, and SC by Definition 9 on PDF p. 13)
% proof_locator: printed p. PDF pp. 16-17 = printed pp. 12:16-12:17 (starts in the last paragraph of PDF p. 16, continues at the top of PDF p. 17, ends with the end-of-proof square just before “Our main theorem now easily follows.”)
% representation: PDF; transcribed from rendered page images (never from the text layer)
% collected_by: carrier rule (the headline theorem was an assembly)
% review_status: two independent reviewers, H2 agreed
% status: complete (statement + author proof verbatim + reason + locators)
%
%% ===========================================================================
%% D3  --  human-proof-corpus candidate transcription
%% ---------------------------------------------------------------------------
%% Source PDF : /disks/sata1/yupeng/human-proof-corpus/source-cache/
%%              registry-sources/drops/t12.pdf        (host Yupeng-orx:443)
%% Paper      : Noel Arteche, Erfan Khaniki, Jan Pich, Rahul Santhanam,
%%              "From Proof Complexity to Circuit Complexity via Interactive
%%               Protocols",
%%              ICALP 2024, LIPIcs Vol. 297, Article No. 12, pp. 12:1-12:20.
%%              DOI 10.4230/LIPIcs.ICALP.2024.12
%% Collected  : Lemma 18 (statement + COMPLETE author proof).
%%              statement : PDF p. 16  = printed p. 12:16
%%              proof     : PDF pp. 16-17 = printed pp. 12:16-12:17
%%              (the proof is split by the page break; the last paragraph of
%%               PDF p. 16 carries the opening case distinction, PDF p. 17
%%               carries the construction and the conclusion, ending in the
%%               end-of-proof square)
%% Context transcribed below for closure (each block clearly delimited):
%%   printed 12:12 : Theorem 7 and Theorem 8 (cross-references; Theorem 8 is
%%                   invoked in the collected proof)
%%   printed 12:13 : sum-check protocol + Definition 9 (SC), which the
%%                   collected statement quantifies over
%%   printed 12:14 : Lemma 10, Theorem 11, Definition 12 (the cited
%%                   quantitative input and the formula being verified)
%%   printed 12:15 : Proposition 13 + its COMPLETE proof (the in-paper step
%%                   that makes the cited input usable in the theory; the
%%                   collected proof consumes its content)
%%   printed 12:16 : Theorem 14, Corollary 15, Lemma 16 + proof,
%%                   Definition 17 (the system named in the statement)
%%   printed 12:17 : Theorem 19 (main theorem) + its proof (shows what the
%%                   collected lemma is the load-bearing step of)
%% Method     : pdftotext -layout was used for LOCATING only.  The text layer
%%              of this PDF is unusable as text of record (it flattens
%%              superscripts, loses blackboard bold, turns the ceiling brackets
%%              of Definition 12 into plain brackets, and mangles \preceq).
%%              Every page transcribed below was rendered with
%%                pdftoppm -f <p> -l <p> -r 150 -png        (whole page)
%%              and every symbol-level decision was taken from crops rendered
%%              with -r 300 or -r 600 -png and read as an image.  The page
%%              images and the crops are kept in b19/pages/ and b19/img/.
%% Typography : reproduced as printed.  \textsc{SC}, \textsc{iEF}, \textsc{EF}
%%              are the small-caps system names of the source.  The source
%%              sets "Proof." in bold; this wrapper sets it bold as well.
%%              Printed slips are kept literally and flagged with
%%              %% [NOTE: ...]; nothing is silently repaired or completed.
%% ILLEGIBLE  : none -- every glyph on the transcribed pages was legible, so
%%              there is no %% [ILLEGIBLE: ...] marker anywhere in this file.
%% ===========================================================================

\documentclass[11pt]{article}
\usepackage[T1]{fontenc}
\usepackage{amsmath,amssymb,amsthm}
\usepackage[margin=1in]{geometry}

\newcommand{\Stwo}{\mathrm{S}^{1}_{2}}
\newcommand{\oneEXP}{\textnormal{1-EXP}}
\newcommand{\SC}{\textsc{SC}}
\newcommand{\iEF}{\textsc{iEF}}
\newcommand{\EF}{\textsc{EF}}
\newcommand{\PV}{\mathrm{PV}}
\newcommand{\GEXP}{\mathrm{G}_{\mathrm{EXP}}}
\newcommand{\Gstar}{\mathrm{G}^{*}_{1}}
\newcommand{\Gone}{\mathrm{G}_{1}}
\newcommand{\Fp}{\mathbb{F}_{p}}
\newcommand{\Fpn}{\mathbb{F}_{p}^{n}}
\newcommand{\N}{\mathbb{N}}
\newcommand{\Size}{\mathrm{Size}}
\newcommand{\HardA}{\mathrm{Hard}^{\mathrm{A}}}
\newcommand{\CFTT}{\mathrm{CorrectFracTT}}
\newcommand{\Snd}{\mathrm{Sound}}
\newcommand{\poly}{\mathrm{poly}}
\newcommand{\RefP}{\mathrm{Ref}}
\newcommand{\ttavg}{\mathrm{tt}^{\mathrm{avg}}}
\newcommand{\iEFtt}[1]{\iEF^{\mathrm{tt}(#1)}}
\newcommand{\ceiling}[1]{\lceil #1 \rceil}
\newcommand{\dotcup}{\mathbin{\dot{\cup}}}
\newcommand{\twohead}{\twoheadrightarrow}

%% statement block: bold run-in label, italic body (as the source sets
%% theorem-like environments); the label is passed as a mandatory argument so
%% that printed bracketed citations such as [39, Thm. 15.3] survive verbatim.
\newenvironment{statement}[1]%
  {\par\medskip\noindent\textbf{#1}\par\nopagebreak\smallskip\noindent\itshape}%
  {\par\medskip}

%% author proof block: bold "Proof." and the source's end-of-proof triangle
\newenvironment{authorproof}%
  {\par\medskip\noindent\textbf{Proof.}\ \ignorespaces}%
  {\unskip\par\nopagebreak\smallskip\noindent\hfill$\blacktriangleleft$\par\medskip}

\begin{document}

%% ###########################################################################
%% ###  CONTEXT  --  printed 12:12 (PDF p. 12): the two correspondence
%% ###  theorems; Theorem 8 is invoked verbatim in the collected proof
%% ###########################################################################

\begin{statement}{$\blacktriangleright$ Theorem 7 (Correspondence of $\Stwo+\oneEXP$ and $\iEF$ [45, Thm.\ 2.1]).}
The proof system $\iEF$ corresponds to $\Stwo+\oneEXP$. That is,
\begin{enumerate}
\item[(i)] the theory $\Stwo+\oneEXP$ proves the soundness of $\iEF$;
\item[(ii)] whenever a $\forall\Pi^{b}_{1}$-sentence $\forall x\varphi(x)$ is provable in
  $\Stwo+\oneEXP$, there are polynomial-size $\iEF$-proofs of the sequence of tautologies
  $\{\|\varphi\|_{n}\}_{n\in\N}$;
\item[(iii)] if $\Stwo+\oneEXP$ proves the soundness of some propositional system $S$, then
  $\iEF \ge S$.
\end{enumerate}
\end{statement}

\begin{statement}{$\blacktriangleright$ Theorem 8 (Correspondence of $\Stwo$ and $\Gstar$ [47]).}
Whenever a $\forall\Sigma^{b}_{1}$-sentence $\forall x\exists y \le t.\varphi(x,y)$ is provable
in $\Stwo$, there are polynomial-size proofs of the sequence of $\Sigma^{q}_{1}$-formulas
obtained by the translation, $\{\|\exists x\varphi(x,y)\|_{n}\}_{n\in\N}$, in $\Gstar$.
\end{statement}

%% [NOTE: Theorem 8 is reproduced exactly as printed, including the free
%%  occurrence of y inside \|\exists x\varphi(x,y)\|_{n}; the sentence being
%%  translated is \forall x\exists y \le t.\varphi(x,y), so the printed
%%  "\exists x" looks like a slip of the source (it should be \exists y).
%%  Nothing is repaired here.]

%% ###########################################################################
%% ###  CONTEXT  --  printed 12:13 (PDF p. 13): the sum-check protocol and
%% ###  Definition 9, i.e. the system SC that Lemma 18 speaks about
%% ###########################################################################

Our proofs rely on a particular interactive protocol, the \emph{Sum-Check Protocol} of Lund,
Fortnow, Karloff, and Nisan [54] for the language of unsatisfiable 3CNFs. Unlike
Merlin-Arthur protocols, this is an interactive protocol running for multiple rounds between
a Prover and a Verifier, before the Verifier makes a decision. We now recall the details of
the protocol.

\medskip\noindent\textbf{The Sum-Check Protocol [54]}

The protocol considers a 3CNF $\varphi(x_{1},\dots,x_{n})$ over $m$ clauses, known to both the
Verifier and the Prover.
\begin{enumerate}
\item The Prover generates a prime number\footnote{The constant $c_{p}$ in the exponent comes
from the formalization of the soundness of the sum-check protocol inside $\Stwo+\oneEXP$ in a
recent work of Khaniki [40]; while we do not need such details in our proofs, we leave it here
to be faithful to the formalization.}
$p \in (2^{2n^{3}+n}, 2^{(2n^{3}+n)^{c_{p}}}]$ together with a Pratt certificate\footnote{A
\emph{Pratt certificate} is a succinct witness for primality checkable in polynomial time [60].
The details are not relevant for our results, but it is important that the Verifier can be
convinced of $p$ being a prime.} on the primality of $p$ and sends them to the Verifier, who
checks for correctness of the certificate, and aborts if incorrect.
\item The Prover and the Verifier arithmetize $\varphi$ into a polynomial
$P_{\varphi}(x_{1},\dots,x_{n})$ of degree at most $3m$ over $\Fp$ in the usual way: a clause
like $(x \vee \neg y \vee z)$ is turned into $1-(1-x)y(1-z)$, and one then takes the product of
all such arithmetized clauses. In this way, for all $x \in \{0,1\}^{n}$,
$\varphi(x) = 1$ if and only if $P_{\varphi}(x) = 1$.
\item The Verifier sets $(a_{1},\dots,a_{n}) := (0,\dots,0)$, $Q_{0}(a_{0}) := 0$ and for
$i \in \{1,\dots,n\}$, the following interaction is carried out:
  \begin{enumerate}
  \item[a.] Leaving $x_{i}$ free, the Prover computes the coefficients of the following
    univariate polynomial over $\Fp$,
    $Q_{i}(x_{i}) := \sum_{x_{i+1}\in\{0,1\}} \cdots \sum_{x_{n}\in\{0,1\}}
     P_{\varphi}(a_{1},\dots,a_{i-1},x_{i},x_{i+1},\dots,x_{n})$
    and sends the $O(m)$ coefficients of $Q_{i}$ to the Verifier.
  \item[b.] The Verifier checks whether $Q_{i}(0)+Q_{i}(1) = Q_{i-1}(a_{i-1})$. If the check
    fails, the Verifier rejects. Otherwise, it samples a random $a_{i} \in \Fp$ and sends it to
    the Prover.
  \item[c.] In the final round, instead of sending $a_{n}$ to the Prover, the Verifier checks
    whether $P_{\varphi}(a_{1},\dots,a_{n}) = Q_{n}(a_{n})$ and accepts or rejects based on this.
  \end{enumerate}
\end{enumerate}

\medskip\noindent\textbf{3\quad Main result}

Our proof exploits the known fact that if $\#\mathbf{P} \subseteq \mathbf{FP}/\mathbf{poly}$,
then $\mathbf{coNP} \subseteq \mathbf{MA}$. Indeed, if $\#\mathbf{P}$ has small circuits one can
provide polynomial-size circuits that simulate the Prover's movements in the Sum-Check protocol
for \textsc{Unsat}, since one can consider the MA proof system in which Arthur receives from
Merlin a circuit claiming to be the circuit that the Prover used to carry out their strategy,
and with the aid of randomness, Arthur can execute this on his own and decide based on the
outcome of this simulation.

Let us make this formal.

\begin{statement}{$\blacktriangleright$ Definition 9 (The $\SC$ proof system).}
Let $V(p,u,\varphi,C,r)$ be the polynomial-time function carrying out the simulation of the
Sum-Check protocol. Namely, $p$ is intended to be a prime in
$(2^{2n^{3}+n}, 2^{(2n^{3}+n)^{c_{p}}}]$, $u$ a Pratt certificate for $p$, $\varphi$ a 3CNF
over $n$ variables, $r$ a string of random bits, and $C$ a multi-output circuit providing the
Prover's responses in the interactions with the Verifier in the Sum-Check protocol.

The Sum-Check Proof System, denoted by $\SC$, is a Merlin-Arthur proof system for proving
3DNF tautologies. An $\SC$ proof of $\varphi$ is a tuple $\langle p,u,C\rangle$ such that $p$ is
indeed a prime in the interval above, correctly certified by the Pratt certificate $u$, and such
that $\Pr_{r\in\Fpn}[V(p,u,\neg\varphi,C,r) = 1] = 1$.
\end{statement}

%% ###########################################################################
%% ###  CONTEXT  --  printed 12:14 (PDF p. 14): Lemma 10 (proof omitted by the
%% ###  authors), Theorem 11 (the cited quantitative soundness bound) and
%% ###  Definition 12 (the formula whose provability Proposition 13 claims)
%% ###########################################################################

The following is just a rephrasing of the fact that $\#\mathbf{P} \subseteq \mathbf{FP}/\mathbf{poly}$
implies $\mathbf{coNP} \subseteq \mathbf{MA}$, in terms of the Merlin-Arthur system $\SC$. We
omit the proof, which can be be found in standard texts (see e.g.\ [5, Thm.\ 8.22]).

%% [NOTE: "can be be found" is printed with the duplicated word in the source;
%%  reproduced literally, not repaired.]

\begin{statement}{$\blacktriangleright$ Lemma 10.}
If $\#\mathbf{P} \subseteq \mathbf{FP}/\mathbf{poly}$, then $\SC$ is polynomially bounded over
3DNF tautologies.
\end{statement}

Our goal is to extend the previous lemma from $\SC$ to the concrete and natural Cook-Reckhow
system Implicit Extended Frege. The idea again is that $\iEF$ (or rather its first-order
counterpart, $\Stwo+\oneEXP$) can prove the soundness of this system and thus simulate it. We
shall then derandomize the $\SC$ protocol inside $\iEF$. Fortunately for us, the soundness of
the Sum-Check protocol was recently proven by Khaniki in the right theory of bounded
arithmetic.

\begin{statement}{$\blacktriangleright$ Theorem 11 (Soundness of the sum-check protocol [39, Thm.\ 15.3]).}
There are constants $c,k \in \N$ such that $\Stwo$ proves the following sentence: for every
$n,\varphi,a,p,u,C$, if it holds that (i) $\varphi$ is a 3CNF in $n$ variables where $n \ge c$,
and, (ii) $\varphi(a) = 1$ and, (iii) $2^{2n^{3}+n} < p \le 2^{(2n^{3}+n)^{c_{p}}}$ and,
(iv) $n^{k} \in \mathrm{Log}\,\mathrm{Log}$, then
\[
  \Pr_{r\in\Fpn}\bigl[V(p,u,\varphi,C,r) = 1\bigr] \;\le\; \frac{n\binom{2^{n}}{3}}{p}.
\]
\end{statement}

We can now formalize the soundness of the $\SC$ proof system from Definition 9. The arguments
that follow are a concrete application of more sophisticated techniques employed by Khaniki
[40, 39], who has studied interactive protocols in the context of defining new jump operators
in proof complexity.

\begin{statement}{$\blacktriangleright$ Definition 12 (The $\Snd_{c}(\SC)$ formula).}
We denote by $\Snd_{c}(\SC)$ the following $\forall\Sigma^{b}_{1}$ sentence, claiming the
soundness of $\SC$: for all $\varphi,a,p,u,C,f$, where $|\varphi| > c$, there is a circuit $D$
of size $\le \ceiling{|f|^{1/4}}$ such that if
\[
  \neg\left(\Pr_{r\in\Fpn}\bigl[V(p,u,\neg\varphi,C,r) = 1\bigr] \preceq^{f}_{\epsilon} \frac{3}{8}\right)
\]
holds, then at least one of the following conditions holds:
\begin{enumerate}
\item[(i)] $|f| \ne |C|^{k_{a}} + k'_{a}$ or,
\item[(ii)] $\CFTT^{1/4}(\ceiling{|f|^{1/4}}, \|f\|, D, f) = 1$ or,
\item[(iii)] $p \notin (2^{2n^{3}+n}, 2^{(2n^{3}+n)^{c_{p}}}]$ or,
\item[(iv)] $\varphi(a) = 1$,
\end{enumerate}
where $k_{a},k'_{a}$ are the constants from Theorem 6 making sure that $\Size$ function works
properly, $\epsilon = 1/16$ and $n$ is the number of variables of $\varphi$. In the definition
of the displayed probability, we assume that $y = p^{n}$ and that the circuit defining the set
of strings accepted by $V$ rejects all $r \ge p^{n}$.
\end{statement}

Note that even if $V$ accepts with probability $1$ on a given input, the approximating
probability from Definition 12 can be significantly smaller because of the difference between
$p^{n}$ and the input-size of the circuit in the input of the $\Size$ function. Another relevant
point is that for each $C,2^{n},e$, the function $\Size(\alpha,C,2^{n},e)$ calls $\alpha$ only
once. In fact, it calls $\alpha(x)$ on an input $x$ which depends only on $|C|,n,|e|$. This is
needed for the formula $\Snd_{c}(\SC)$ to be well-defined.

It now suffices to verify that the encoding of the soundness of $\SC$ is indeed provable in
$\Stwo+\oneEXP$.

%% ###########################################################################
%% ###  CONTEXT  --  printed 12:15 (PDF p. 15): Proposition 13 with its
%% ###  COMPLETE proof.  This is the in-paper step that turns the cited
%% ###  Theorem 11 into a statement usable in the theory; the proof of the
%% ###  COLLECTED Lemma 18 consumes it (through "by Theorem 11 and
%% ###  Corollary 15" on printed 12:17).
%% ###########################################################################

\begin{statement}{$\blacktriangleright$ Proposition 13 (Soundness of $\SC$ inside $\Stwo+\oneEXP$).}
There is a constant $c \in \N$ such that $\Stwo+\oneEXP \vdash \Snd_{c}(\SC)$.
\end{statement}

\begin{authorproof}
Let $c \in \N$ be a big enough constant that can be computed from the rest of the argument and
\[
  \Snd_{c}(\SC) := \forall\varphi,a,p,u,C,f\exists D\Phi(\varphi,a,p,u,C,f,D)
\]
the soundness formula in Definition 12 above. Let $\varphi$ be a 3DNF in $n$ variables such that
$|\varphi| > c$, and consider $a,p,u,C,f$. Then the following cases can happen:
\begin{enumerate}
\item[(a)] If $|f| \ne |C|^{k_{a}}+k'_{a}$ or $p \notin (2^{2n^{3}+n}, 2^{(2n^{3}+n)^{c_{p}}}]$,
  then $\Phi(\varphi,a,p,u,C,f,0)$ is trivially true.
\item[(b)] If there is a circuit $D$ of size $\le \ceiling{|f|^{1/4}}$ such that
  $\CFTT^{1/4}(\ceiling{|f|^{1/4}},\|f\|,D,f) = 1$, then $\Phi(\varphi,a,p,u,C,f,D)$ is trivially
  true.
\item[(c)] If the previous cases do not happen and moreover
\[
  \neg\left(\Pr_{r\in\Fpn}\bigl[V(p,u,\neg\varphi,C,r) = 1\bigr] \preceq^{f}_{\epsilon} \frac{3}{8}\right)
\]
  holds, then we have that $8 \cdot \Size(f,C^{*},2^{m},e) > 3p^{n}$, where $m$ is the smallest
  integer such that $2^{m} \ge p^{n}$, $\epsilon := |e|^{-1}$ and
  $C^{*}(r) := V(p,u,\neg\varphi,C,r)$. By the assumption $\HardA_{1/4}(f)$ holds and by the
  fact that we are over $\Stwo$ and we can use $f$ as a parameter in polynomial induction for
  $\Sigma^{b}_{1}$ formulas, we can do approximate counting using Theorem 6. Hence there is a
  circuit $G$ and some $v \le \poly(m\epsilon^{-1}|C^{*}|)$ such that
\[
  G : v \times (X_{C^{*}} \dotcup \epsilon 2^{m}) \twohead v \times \Size(f,C^{*},2^{m},e).
\]
  As we work in $\Stwo+\oneEXP$ and $G$ is surjective, we can find a subset
  $A \subseteq v \times (X_{C^{*}} \dotcup \epsilon 2^{m})$ such that $G$ restricted to $A$ is a
  one-to-one function from $A$ to $v \times \Size(f,C^{*},2^{m},e)$. Now we can apply exact
  counting (as we have $\oneEXP$) and show that
\[
  \Size(f,C^{*},2^{m},e) \le |X_{C^{*}}| + \epsilon 2^{m}.
\]
  By the fact that $8 \cdot \Size(f,C^{*},2^{m},e) > 3p^{n} > 3 \cdot 2^{m}/2$, we have
  $2^{m}/8 < |X_{C^{*}}|$. Now if $\varphi(a) = 0$, by Theorem 11 we get
\[
  \Pr_{r\in\Fpn}\bigl[V(p,u,\neg\varphi,\pi,r) = 1\bigr] \;\le\; \frac{n\binom{2^{n}}{3}}{p}.
\]
%% [NOTE: the argument position here is printed "\pi", while the object in play
%%  is the circuit C (and C^{*} everywhere else in this proof).  Reproduced as
%%  printed; the same slip occurs in the next display as well.]
  Note that $|\varphi| > c$ which implies that $n$ is big enough and as $p > 2^{2n^{3}+n}$ we get
  that $n\binom{2^{n}}{3}/p \le 1/8$, which implies
\[
  \Pr_{r\in\Fpn}\bigl[V(p,u,\neg\varphi,\pi,r) = 1\bigr] \;\le\; \frac{1}{8}.
\]
  As $C^{*}$ rejects all $r \ge p^{n}$, this implies that $|X_{C^{*}}| \le 2^{m}/8$ which leads
  to a contradiction, so $\varphi(a) = 1$.
\end{enumerate}
\end{authorproof}

The main technical issue now is that $\Snd_{c}(\SC)$ is a $\forall\Sigma^{b}_{1}$ sentence that
does not translate into a propositional formula that $\iEF$ can reason about. Instead, we shall
work on a quantified propositional system. For this to make sense we need to know the quantified
propositional proof system associated with $\Stwo+\oneEXP$.

We invoke the following known \textsc{tfnp} characterization of the $\Sigma^{b}_{1}$ consequences
of $\Stwo+\oneEXP$, which identifies a ``complete'' $\Sigma^{b}_{1}$ sentence $\Psi$ such that any
other $\Sigma^{b}_{1}$ consequence of $\Stwo+\oneEXP$ reduces to it in $\Gstar$.

%% ###########################################################################
%% ###  CONTEXT  --  printed 12:16 (PDF p. 16): Theorem 14, Corollary 15
%% ###  (the authors state that the proofs of its items are only in the full
%% ###  version of the paper), Lemma 16 with its proof, and Definition 17,
%% ###  which defines the system named in the COLLECTED statement.
%% ###########################################################################

\begin{statement}{$\blacktriangleright$ Theorem 14 ([42, 41, 46, 9]).}
There is a $\forall\Sigma^{b}_{1}$ sentence $\Psi := \forall x\exists y\psi(x,y)$ (the bound on
$y$ is implicit in $\psi$) such that the following statements are true:
\begin{enumerate}
\item[(i)] $\Stwo+\oneEXP \vdash \forall x\exists y\psi(x,y)$;
\item[(ii)] for any $\forall\Sigma^{b}_{1}$ sentence $\forall x\exists y\alpha(x,y)$ such that
  $\Stwo+\oneEXP \vdash \forall x\exists y\alpha(x,y)$, there are $\PV$ functions $f$ and $g$ such
  that $\Stwo \vdash \forall x,y(\psi(f(x),y) \rightarrow \alpha(x,g(x,y)))$.
\end{enumerate}
\end{statement}

In what follows, we shall work with Gentzen's system $\mathrm{G}$ extended with the propositional
translation of the sentence $\Psi$ in the theorem above. We denote this system by
$\GEXP := \Gstar + \|\Psi\|$ and use the following key properties about it. The full version of
the paper contains explicit proofs of each of these items, but we shall omit these here.

\begin{statement}{$\blacktriangleright$ Corollary 15.}
The following statements about $\GEXP$ hold:
\begin{enumerate}
\item[(i)] $\Stwo+\oneEXP \vdash \Sigma^{q}_{1}$-$\RefP(\GEXP)$, i.e.\ the reflection principle
  for $\GEXP$ and $\Sigma^{q}_{1}$ formulas is provable in $\Stwo+\oneEXP$;
\item[(ii)] for every $\forall\Sigma^{b}_{1}$-sentence $\forall x\exists y\alpha(x,y)$, if
  $\Stwo+\oneEXP \vdash \forall x\exists y\alpha(x,y)$, then there are polynomial-size
  $\GEXP$-proofs of the sequence of $\Sigma^{q}_{1}$-tautologies
  $\{\|\exists y\alpha(y)\|_{n}\}_{n\in\N}$;
\item[(iii)] if $\Stwo+\oneEXP$ proves the soundness of a propositional proof system $S$, then
  $\GEXP \ge S$.
\end{enumerate}
\end{statement}

Let us observe that $\GEXP$ is in fact equivalent to $\iEF$.

\begin{statement}{$\blacktriangleright$ Lemma 16.}
The proof systems $\iEF$, $\EF+\RefP_{\iEF}$ and $\GEXP$ are polynomially equivalent over
propositional tautologies.
\end{statement}

\begin{authorproof}
By item (iii) of Corollary 15 and item (iii) Theorem 7, $\iEF$ and $\GEXP$ polynomially simulate
each other. As mentioned in Section 2.1.1, $\EF+\RefP_{\iEF} \ge \iEF$. It is also easy to see
that $\Stwo+\oneEXP$ proves the soundness of $\EF+\RefP_{\iEF}$, which by item (iii) of Theorem 7
gives us $\iEF \ge \EF+\RefP_{\iEF}$.
\end{authorproof}

We are now ready to define the extension of $\iEF$ for which our main theorem holds. Recall that
the propositional formulas $\ttavg_{1/4}(h_{n},2^{n/4})$ were defined in Section 2.2.2 and state
the average-case hardness of a Boolean function $h_{n}$ represented as a truth table.

\begin{statement}{$\blacktriangleright$ Definition 17 (The systems $\iEF^{\mathrm{tt}}$).}
Let $h = \{h_{n}\}_{n\in\N}$ be some family of Boolean functions, and let $n_{0} \in \N$. We
denote by $\iEFtt{h,n_{0}} := \GEXP + \{\ttavg_{1/4}(h_{n},2^{n/4})\}_{n\ge n_{0}}$ the system
that extends $\GEXP$ by the axioms claiming the hardness of $h_{n}$, for $n \ge n_{0}$.
\end{statement}

Note that $\iEFtt{h,n_{0}}$ is a family of proof systems, parameterized by a Boolean function
family $h$ and some threshold parameter $n_{0}$. Observe that depending on the choice of $h$ and
$n_{0}$, the system $\iEFtt{h,n_{0}}$ may not be a Cook-Reckhow system: if $h$ is not a hard
function, or $n_{0}$ is not large enough, we will be adding axioms which are not tautologies; and
even if $h$ is hard and $n_{0}$ is large enough, the system may require advice to verify the
proofs. As we shall see, however, these degenerate instantiations of $\iEFtt{h,n_{0}}$ are not a
problem.

What is more important, the systems $\iEFtt{h,n_{0}}$, regardless of their consistency, always
simulate $\SC$.

%% ###########################################################################
%% ###                                                                    ####
%% ###   PRIMARY COLLECTED UNIT  --  Lemma 18                              ####
%% ###   printed 12:16 (statement + first paragraph of the proof)          ####
%% ###   printed 12:17 (rest of the proof, up to the end-of-proof square)  ####
%% ###                                                                    ####
%% ###########################################################################

\begin{statement}{$\blacktriangleright$ Lemma 18.}
Let $h$ be family of Boolean functions and let $n_{0} \in \N$. The system $\iEFtt{h,n_{0}}$
polynomially simulates $\SC$ over 3DNF tautologies.
\end{statement}

%% [NOTE: "Let $h$ be family of Boolean functions" is printed without the
%%  article in the source; reproduced literally.]

\begin{authorproof}
If the system $\iEFtt{h,n_{0}}$ is unsound because the added axioms are not tautologies, then the
system is polynomially bounded and simulates every other proof system. So suppose the added
axioms are indeed tautologies, meaning that the function $h$ is indeed hard on average.

Let $\varphi_{1}$ be a 3DNF in $n_{1}$ variables and $\langle p_{1},u_{1},C_{1}\rangle$ be a
$\SC$-proof of $\varphi_{1}$. This means
\[
  2^{2n_{1}^{3}+n_{1}} < p_{1} \le 2^{(2n_{1}^{3}+n_{1})^{c_{p}}} \;\wedge\;
  \Pr_{r\in\mathbb{F}_{p_{1}}^{n_{1}}}\bigl[V(p_{1},u_{1},\neg\varphi_{1},C_{1},r) = 1\bigr] = 1.
\]
Note that by Theorem 11 and Corollary 15, there are $\PV$ functions $l,g$ such that
\[
  \Stwo \vdash \forall\varphi,a,p,u,C,f\,
  \bigl(\psi(l(\varphi,a,p,u,C,f),y) \rightarrow
  \Phi(\varphi,a,p,u,C,f,g(\varphi,a,p,u,C,f,y))\bigr),
\]
where $\Snd_{c}(\SC) := \forall\varphi,a,p,u,C,f\exists D\Phi(\varphi,a,p,u,C,f,D)$. Let
$s := |\langle p,u,C\rangle|$. Then by Theorem 8 there is a $s^{O(1)}$-size $\Gstar$-proof of
\[
  \bigl\|\forall\varphi,a,p,u,C,f\,\bigl(\psi(l(\varphi,a,p,u,C,f),y) \rightarrow
  \Phi(\varphi,a,p,u,C,f,g(\varphi,a,p,u,C,f,y))\bigr)\bigr\|_{s'},
\]
where $s' := \poly(s)$.

%% [NOTE: the source prints "Let $s := |\langle p,u,C\rangle|$" with
%%  unsubscripted p, u, C, although the objects in play are p_1, u_1, C_1
%%  (they are subscripted again in the very next sentence).  Reproduced as
%%  printed.]

Let us rewrite the previous quantified propositional formula as
$\|\Psi'\| \rightarrow \|\Phi'\|$ with the right range of parameters such that
$p_{1},u_{1},\varphi_{1},C_{1}$ are substituted in the formula in their corresponding places.
Now we take the substitution instance $\ttavg_{1/4}(h_{n'},2^{n'/4})$ where
$|h_{n'}| := |C_{1}|^{k_{a}} + k'_{a}$ and we substitute $h_{n'}$ to the variables corresponding
to $f$ and therefore the disjunct which corresponds to $\CFTT$ disappears from $\|\Phi'\|$ when
we apply the rules of $\Gstar$. Moreover, it is not hard to verify that after the substitutions
every other disjunct which corresponds to subformulas of $\Snd_{c}(\SC)$ from Definition 12
disappears except $\varphi_{1}$. So what we have is $\Gstar$-proof of
$\|\Psi''\|(\bar{x},\bar{y}) \rightarrow \varphi_{1}(\bar{x})$ ($\bar{x}$ and $\bar{y}$ are
disjoint variables) where $\|\Psi''\|$ is a substitution instance of $\|\Psi'\|$. Since we are
working in $\GEXP$, we have the substitution instance $\exists\bar{y}\|\Psi''\|(\bar{x},\bar{y})$
and therefore using the rules of $\Gstar$ we get a short $\GEXP$-proof of $\varphi_{1}(\bar{x})$.
\end{authorproof}

%% ###########################################################################
%% ###  CONTEXT  --  printed 12:17 (PDF p. 17): Theorem 19, the main theorem,
%% ###  whose proof is the one-line reduction to the collected Lemma 18 and
%% ###  to Lemma 10 (whose proof the authors omit)
%% ###########################################################################

Our main theorem now easily follows.

\begin{statement}{$\blacktriangleright$ Theorem 19 (Main theorem).}
Let $h$ be a family of Boolean functions and let $n_{0} \in \N$. If the system $\iEFtt{h,n_{0}}$
is not polynomially bounded, then $\#\mathbf{P} \not\subseteq \mathbf{FP}/\mathbf{poly}$.
\end{statement}

\begin{authorproof}
By Lemma 18 above, for every choice of $h$ and $n_{0}$, the system $\iEFtt{h,n_{0}}$
polynomially simulates $\SC$, so if $\iEFtt{h,n_{0}}$ is not polynomially bounded, then $\SC$ is
not either. Then, by the contrapositive of Lemma 10,
$\#\mathbf{P} \not\subseteq \mathbf{FP}/\mathbf{poly}$.
\end{authorproof}

As discussed, depending on the choice of $h$ and $n_{0}$, the system $\iEFtt{h,n_{0}}$ may not be
sound and thus possibly not Cook-Reckhow. However, for any fixed choice of a uniform candidate
hard function, the system is concrete and exhibits the desired connection that proof complexity
lower bounds for it imply strong circuit lower bounds. In particular, if there exist functions
in $\mathbf{NE} \cap \mathbf{coNE}$ average-case hard for subexponential-size circuits, then we
recover the version of the theorem presented in the introduction (Theorem 1).

\end{document}
